\documentclass[12pt,a4paper]{article}

\usepackage[margin=2.5cm]{geometry}
\usepackage{amsmath,amssymb,amsthm}
\usepackage{physics}
\usepackage{graphicx}
\usepackage{hyperref}
\usepackage{bm}

\title{Compatibility of the Schr\"odinger Framework with the Projective Dynamic Logo (PDL)}
\author{C.~Laubscher}
\date{\today}

\begin{document}

\maketitle

\begin{abstract}
We examine the consistency between the standard nonrelativistic Schr\"odinger description of the hydrogen atom and the Projective Dynamic Logo (PDL), a relational framework in which elementary entities are represented as minimal logical closures on signed graphs. Without attempting to derive the Schr\"odinger equation from the axioms of PDL, we demonstrate that, once Schr\"odinger dynamics for an effective degree of freedom is postulated, both the Coulomb coupling and the fine-structure constant can be reinterpreted as emergent quantities arising from the combinatorial structure of the proton and the elementary $(4,6)$ block associated with the electron. Within controlled uncertainties, the PDL-based expression for the fine-structure constant reproduces its empirical value, and the resulting effective Coulomb potential yields the correct order of magnitude for the hydrogen ground-state binding energy and the Bohr radius. The main contribution of this work lies in providing a structural origin and interpretation of the coupling, rather than proposing any modification of the established Schr\"odinger phenomenology.
\end{abstract}

\section{Introduction}

\subsection{Background and motivation}

Nonrelativistic quantum mechanics yields an exceptionally accurate description of the hydrogen atom. In its conventional formulation, the electron is described by a wave function \(\psi(\mathbf{r},t)\) that evolves according to the Schrödinger equation in the presence of a Coulomb potential generated by a pointlike proton,
\[
  i\hbar \frac{\partial \psi}{\partial t}(\mathbf{r},t)
  =
  \left[
    -\frac{\hbar^2}{2m_e}\Delta
    - \frac{e^2}{4\pi\varepsilon_0}\,\frac{1}{r}
  \right]
  \psi(\mathbf{r},t),
\]
where \(m_e\) denotes the electron mass and \(e^2/(4\pi\varepsilon_0)\) is regarded as a fixed coupling constant. The associated time-independent eigenvalue problem reproduces the well-known hydrogenic energy spectrum and the Rydberg series, with the fine-structure constant \(\alpha = e^2/(4\pi\varepsilon_0 \hbar c)\) appearing as a dimensionless parameter that quantifies the strength of the electromagnetic interaction.

Within this standard theoretical context, the quantities \(\hbar\), \(m_e\), \(m_p\), and \(\alpha\) are regarded as empirically determined input parameters. Although their numerical values can be measured with very high precision, the framework itself does not provide a structural or mechanistic explanation for the specific magnitudes they assume. From a foundational perspective, it is therefore natural to ask whether at least a subset of these constants might be derivable from, or attributable to, a more primitive combinatorial or relational substrate, rather than being posited as irreducible attributes of pre-given fields or particles.

The Projective Dynamic Logo (PDL) furnishes a candidate formalism for such a reconstruction.\cite{Laubscher2026IntroPDL,Laubscher2026PDL_English,Laubscher2026LDP_French}In PDL, physical reality is modelled as a network of minimal logical closures represented by finite signed graphs, whose edges encode binary agreement/disagreement relations and whose configurations evolve according to a coherence-optimisation principle, without any prior assumption of an underlying spacetime manifold or material substratum.The framework is based on four axioms: minimal binary pulsation, triangular coherence, minimal completeness, and logical optimisation.Subject to these axioms, the first admissible elementary stationary configuration is identified as the complete signed graph on four vertices and six edges, denoted \((4,6)\). For this graph, there exists a sign assignment that renders all triangular subgraphs coherent, together with a global sign inversion that leaves this coherence invariant.This minimally closed configuration is naturally interpreted as corresponding to an electron at rest.

At a higher level of structural complexity, the same axioms single out a hierarchical combinatorial architecture for the proton, comprising three local stationary domains (valence cores) embedded within a dense relational background (relational sea). In this formulation, the valence sector is specified by integer parameters \(n_u = 24\) and \(n_d = 28\) for up- and down-type cores, respectively, yielding a total of
\[
  r_{\text{val}} = 930
\]
stationary relations, whereas the relational sea contributes
\[
  R_{\text{sea}} = 10\,087
\]
dynamical links. The resulting total relational content is therefore
\[
  R_{\text{tot}} = r_{\text{val}} + R_{\text{sea}} = 11\,017.
\]
Accordingly, approximately \(9\%\) of the internal relations are stationary and \(91\%\) are dynamical. This distribution is in qualitative agreement with the standard QCD-based picture, in which only a small fraction of the proton mass originates from valence degrees of freedom.\cite{Laubscher2026PDL_English,Laubscher2026LDP_French}

Within this relational framework, several dimensionless and scale-setting constants can be reformulated as ratios of internal coherence. The reduced Planck constant \(\hbar\) is interpreted as the minimal quantum of action associated with a complete coherence cycle of the elementary \((4,6)\) structure. The fine-structure constant \(\alpha\) admits a reinterpretation as a surface-to-volume coherence fraction, quantifying the ratio between the global coherence budget of the proton and the portion effectively projected onto a finite active surface. Boltzmann’s constant \(k_B\) and an effective gravitational coupling can likewise be expressed in terms of the same integers, the proton–electron mass ratio, and a small coherence-leakage parameter, although these constructions remain partly conjectural.\cite{Laubscher2026PDL_English,Laubscher2026LDP_French}

In the present analysis, we do not attempt to replace the Schr\"odinger description of hydrogen by a fully relational dynamical theory. Instead, we treat the standard nonrelativistic Schr\"odinger framework as an effective description of an electronic degree of freedom, and investigate whether the Coulomb coupling entering this framework can be reconstructed from the internal combinatorial structure of the proton together with the minimal \((4,6)\) block associated with the electron. The central question is whether the conventional Coulomb potential can be preserved in its usual form, while assigning its coupling coefficient a structural interpretation derived from PDL.

\subsection{Scope and objectives}

The purpose of this article is to formulate and analyse a concrete form of compatibility between the standard Schr\"odinger description of the hydrogen atom and the PDL relational framework. More specifically, the analysis proceeds in three stages.

First, we review the PDL construction of the electron and the proton. An electron at rest is represented by the minimal stationary closure \((4,6)\), whose internal binary pulsation is interpreted as the logical counterpart of the electron Compton pulsation, and whose coherence cost per cycle is quantified by \(\hbar\). The proton is modelled as a two-level hierarchical architecture composed of three valence cores of sizes \(n_u = 24\) and \(n_d = 28\), a dense dynamical sea containing \(R_{\text{sea}} = 10\,087\) relations, and a finite active surface mediating couplings to external \((4,6)\)-type structures. This architecture fixes a valence content \(r_{\text{val}} = 930\), a total relational content \(R_{\text{tot}} = 11\,017\), and stationary/dynamical fractions of approximately \(9\%\) and \(91\%\), respectively.\cite{Laubscher2026PDL_English,Laubscher2026LDP_French}

Second, we refine the notion of active surface and derive from it an expression for the fine-structure constant. The effective number of surface relations is estimated as
\[
  R_{\text{surf}} = \varphi\,\frac{r_{\text{val}}}{3},
\]
where \(\varphi = (1+\sqrt{5})/2\) is proposed as a candidate relational optimisation factor and \(r_{\text{val}}/3\) denotes an average coherence core per valence quark. The PDL expression for the inverse fine-structure constant then takes the form
\[
  \alpha_{\text{PDL}}^{-1}
  =
  \frac{\mu}{R_{\text{surf}}^2}
  =
  \frac{\mu}{\varphi^2}\left(\frac{3}{r_{\text{val}}}\right)^2,
\]
where \(\mu = m_p/m_e\) is the proton–electron mass ratio. For \(r_{\text{val}} = 930\), \(\mu \simeq 1836.15\) and \(R_{\text{surf}} \simeq 501.59\), this relation yields \(\alpha_{\text{PDL}}^{-1} \simeq 137.02\), which agrees with the empirical value \(\alpha^{-1} \simeq 137.036\) at the level of a few \(10^{-4}\) in relative terms, without the introduction of any additional continuous free parameter.\cite{Laubscher2026PDL_English,Laubscher2026LDP_French}

Third, we substitute \(\alpha_{\text{PDL}}\) into the standard hydrogenic formulas. Retaining the usual Schr\"odinger equation with Coulomb potential, but replacing \(\alpha\) by \(\alpha_{\text{PDL}}\), we obtain an effective potential
\[
  V_{\text{eff}}(r) = -\alpha_{\text{PDL}}\hbar c\,\frac{1}{r},
\]
which leads to bound-state energies
\[
  E_n = -\frac{1}{2} m_e c^2 \frac{\alpha_{\text{PDL}}^2}{n^2}
\]
and to a modified Bohr radius
\[
  a_0^{\text{PDL}} = \frac{\hbar}{m_e c\,\alpha_{\text{PDL}}}.
\]
The deviations from the standard predictions are of the same order of magnitude as the relative difference between \(\alpha_{\text{PDL}}\) and \(\alpha\), and remain compatible with the observed hydrogen binding energy at the Bohr scale. Within the approximations employed, this indicates that the Schr\"odinger phenomenology can be preserved while the Coulomb coupling acquires a structural origin in terms of surface coherence in the PDL architecture.

The remainder of the paper is organised as follows.
Section~2 briefly recalls the PDL axioms and the identification of the minimal \((4,6)\) structure with an electron at rest. Section~3 summarises the hierarchical architecture of the proton, the construction of its active surface, and the resulting expression for \(\alpha_{\text{PDL}}\). Section~4 develops the effective Coulomb potential and the hydrogenic spectrum obtained by inserting \(\alpha_{\text{PDL}}\) into the standard Schr\"odinger framework. Section~5 examines the conceptual implications and limitations of this compatibility result.

\section{PDL framework and elementary structures}

\subsection{Relational axioms and signed graphs}

Within the PDL framework, physical reality is represented as a network of minimal logical closures formalised by finite signed graphs.\cite{Laubscher2026IntroPDL,Laubscher2026PDL_English,Laubscher2026LDP_French} A logical configuration is specified by a finite graph \(G = (V,E)\), where the vertex set \(V\) encodes informational entities and the edge set \(E\) encodes binary agreement/disagreement relations between them. Each edge \((i,j)\in E\) carries an associated sign \(s_{ij}\in\{+1,-1\}\), on which an elementary binary pulsation acts as a discrete dynamical process.

Triangles, i.e. cycles of length three, constitute the elementary motifs of logical closure. For three vertices \(i,j,k\in V\), the corresponding triangle is present if the three edges \((i,j),(j,k),(i,k)\) all belong to \(E\). The triangle is said to be \emph{coherent} if the product of the associated edge signs satisfies
\[
  s_{ij}\,s_{jk}\,s_{ik} = +1,
\]
and \emph{violated} otherwise, either because at least one of the three edges is absent from \(E\) or because the product of the signs equals \(-1\). Coherent triangles implement minimal logical closure in the sense that the reference structure required to maintain the distinctions between the three vertices is entirely contained within the cycle itself.\cite{Laubscher2026IntroPDL,Laubscher2026PDL_English}

To quantify deviations from triangular closure, we introduce the fraction of violated triangles
\[
  \eta(G)
  =
  \frac{N_{\text{viol}}(G)}{\binom{n}{3}},
\]
where \(n = |V|\) and \(N_{\text{viol}}(G)\) denotes the number of violated triangles in the signed configuration \((G,s)\). By convention, we set \(\eta(G)=0\) whenever \(n<3\). Configurations with small values of \(\eta(G)\) exhibit a high degree of triangular coherence and thus a pronounced internal stability. The quantity \(\eta(G)\) plays a central role in the logical optimisation principle that selects dynamically sustainable structures.\cite{Laubscher2026IntroPDL,Laubscher2026PDL_English}

In addition to coherence, two further ingredients are required. First, the \emph{relational density} is defined as
\[
  \rho(G) = \frac{r}{r_{\max}(n)},
  \qquad
  r = |E|,
  \qquad
  r_{\max}(n) = \frac{n(n-1)}{2},
\]
such that \(\rho(G)=1\) if and only if \(G\) is a complete graph. Second, we introduce a \emph{minimality indicator} \(m(G)\), which takes the value \(1\) if the configuration realises a closed structure that does not require any strictly richer superstructure to achieve closure, and \(0\) otherwise. The three quantities \(\eta\), \(\rho\), and \(m\) are then combined into a logical optimisation functional that discriminates between merely admissible configurations and dynamically persistent ones.\cite{Laubscher2026IntroPDL,Laubscher2026PDL_English}

We model this optimisation functional as
\[
  F(G)
  =
  w_{\text{coh}}\,f_{\text{coh}}\!\bigl(\eta(G)\bigr)
  +
  w_{\text{conn}}\,f_{\text{conn}}\!\bigl(\rho(G)\bigr)
  +
  w_{\min}\,m(G),
\]
where \(w_{\text{coh}},w_{\text{conn}},w_{\min}>0\) are fixed weights and \(f_{\text{coh}},f_{\text{conn}}\) are monotone functions chosen such that:
(i) \(F(G)\) strictly decreases as \(\eta(G)\) increases (configurations with fewer coherence violations are favoured),
(ii) \(F(G)\) increases with \(\rho(G)\) over the relevant range (higher relational density up to closure is preferred),
(iii) \(F(G)\) is larger when \(m(G)=1\) (elementary closed structures are preferred over non-minimal ones).
In the simplest implementation, one may choose \(f_{\text{coh}}(\eta)=1-\eta\) and \(f_{\text{conn}}(\rho)=\rho\), together with \(w_{\text{coh}}=w_{\text{conn}}=w_{\min}=1\). This choice is sufficient to single out minimal closed configurations, such as the \((4,6)\) structure at the fundamental level.\cite{Laubscher2026IntroPDL,Laubscher2026PDL_English}

\subsection{The minimal \texorpdfstring{\((4,6)\)}{(4,6)} structure}

Within this framework, an elementary stationary structure is defined as a configuration \((G,s)\) such that:
(i) \(G\) is connected and complete;
(ii) there exists a sign assignment \(s\) for which all triangular subgraphs are coherent, i.e. \(\eta(G)=0\);
(iii) the union of coherent triangles forms a connected subgraph containing at least two triangles sharing a common edge;
(iv) all vertices are structurally equivalent, so that no internal hierarchy is present;
(v) the discrete sign dynamics admits a binary pulsation \(s\mapsto -s\) that preserves triangular coherence.\cite{Laubscher2026IntroPDL,Laubscher2026PDL_English}

An elementary stationary structure is said to be minimal if no proper subgraph of \(G\) satisfies all of these conditions.
A systematic exclusion of non-admissible configurations shows that no graph with fewer than four vertices can realise such a structure: for one or two vertices, no triangle can be formed, and for three vertices there is at most a single triangle, which is insufficient to sustain a network of overlapping coherent cycles.
The first admissible instance arises for \(n=4\), where the complete graph has
\[
  r_4 = \frac{4\times 3}{2} = 6
\]
edges and is uniquely characterised, up to relabelling, by the pair \((4,6)\).\cite{Laubscher2026IntroPDL,Laubscher2026PDL_English}
For this graph, there exists a sign assignment such that all triangular faces are coherent, and a global sign inversion \(s\mapsto -s\) generates a binary pulsation that leaves the coherence pattern invariant.
The \((4,6)\) configuration thus provides the first explicit realisation of the minimal completeness requirement and functions as the fundamental building block of the PDL architecture.

From a physical perspective, the \((4,6)\) structure is interpreted as the logical prototype of an electron at rest.\cite{Laubscher2026PDL_English,Laubscher2026LDP_French}
By construction, it realises a closed network consisting of four informational entities linked by six signed relations, such that all triangular subgraphs are mutually coherent and a global sign inversion induces an intrinsic logical 2-cycle.
This internal pulsation is independent of any external time parameter and is identified with the abstract analogue of the electron’s Compton pulsation.
The reduced Planck constant \(\hbar\) is then interpreted as the minimal quantum of action associated with a single coherence cycle of this internal dynamics, and the electron rest energy
\[
  E_{0,e} = \hbar\,\omega_{\text{C},e}
\]
is reinterpreted as the coherence cost required to sustain the elementary stationary regime corresponding to the \((4,6)\) closure.

\subsection{Energy as coherence cost}

Within the PDL framework, energy is not postulated as a primitive quantity attributed to entities defined on a pre-given substrate.
Rather, it is introduced as the quantitative measure of the coherence cost required to maintain a given stationary regime.\cite{Laubscher2026IntroPDL,Laubscher2026PDL_English}
For a structure endowed with an intrinsic pulsation of angular frequency \(\omega\), the standard relation
\[
  E = \hbar\,\omega
\]
is reinterpreted as expressing that each coherence cycle is associated with an irreducible quantum of action, determined by \(\hbar\).
The minimal \((4,6)\) configuration thereby establishes the fundamental discreteness of action: its internal pulsation realises the first admissible logical oscillation, and \(\hbar\) quantifies the minimal increment by which internal coherence can be closed upon itself.

Once stationary regimes are endowed with distinct coherence costs, energy acquires a strictly relational and structural status.
One can then distinguish multiple stationary regimes coexisting within a single relational architecture, each characterised by its own intrinsic pulsation and associated action per cycle, without the need to assume any underlying spacetime metric.
In this perspective, the electron at rest emerges as the first non-trivial, self-sustaining regime of existence permitted by the axioms, and its rest energy quantifies the minimal cost necessary to sustain the elementary \((4,6)\) closure over time.\cite{Laubscher2026PDL_English,Laubscher2026LDP_French}

\section{Relational Architecture of the Proton}

\subsection{Valence Cores and Relational Sea}

Whereas the electron is described by a single elementary stationary closure of type \((4,6)\), the proton requires a substantially more intricate composite structure. From the perspective of quantum chromodynamics (QCD), the proton mass is known to arise predominantly from gluonic fields and sea-quark contributions, with the bare masses of the valence quarks accounting for only a minor fraction of the total energy. Within the PDL framework, this empirical complexity is encoded in a two-level hierarchical organisation that combines three local stationary domains (valence cores) with a dynamically fluctuating relational sea.\cite{Laubscher2026PDL_English,Laubscher2026LDP_French}

Local stationary regimes are postulated to satisfy a minimal completeness requirement: the number \(n_q\) of elementary entities constituting a valence quark must be a multiple of four, so that its internal structure can be represented as a superposition of logical tetrads. For a generic quark, this leads to the constraint
\[
  n_q \in \{4,8,12,16,20,24,28,\dots\},
  \qquad
  r_q = \frac{n_q(n_q-1)}{2},
\]
where \(r_q\) denotes the total number of internal relations. In addition, a minimal degree of structural asymmetry is required to generate a non-vanishing net electric charge: two valence quarks (of up type) must share the same internal organisation, whereas the third (of down type) must realise a slightly different configuration. Under these conditions, the pair \((n_u,n_d) = (24,28)\) is identified as the lowest-order configuration that simultaneously ensures (i) internal completeness of each quark, (ii) minimal asymmetry sufficient to produce a total charge \(+1\), and (iii) a maximal value of the optimisation functional \(F\) within the class of configurations examined.\cite{Laubscher2026LDP_French,Laubscher2026PDL_English}

For this choice, the numbers of stationary relations associated with the up- and down-type valence cores are
\[
  r_u = \frac{24\times 23}{2} = 276,
  \qquad
  r_d = \frac{28\times 27}{2} = 378,
\]
so that the total number of stationary valence relations becomes
\[
  r_{\text{val}} = 2r_u + r_d = 2\times 276 + 378 = 930.
\]
The relational sea is represented as a sparse graph with minimal connectivity, characterised by the simple counting rule \(R_{\text{sea}} = 2\,n_{\text{sea}}\). For the proton, this yields
\[
  R_{\text{sea}} = 10\,087,
\]
which is interpreted as the dynamical, non-stationary contribution to the global coherence. The total internal relational content of the proton is therefore
\[
  R_{\text{tot}} = r_{\text{val}} + R_{\text{sea}} = 930 + 10\,087 = 11\,017.
\]
From this, one obtains the stationary and dynamical fractions
\[
  f_{\text{stat}} = \frac{r_{\text{val}}}{R_{\text{tot}}} \simeq 0.09,
  \qquad
  f_{\text{dyn}} = \frac{R_{\text{sea}}}{R_{\text{tot}}} \simeq 0.91,
\]
which are qualitatively consistent with the QCD-based picture, according to which approximately \(90\%\) of the proton’s energy derives from internal dynamical degrees of freedom rather than from the bare rest masses of the valence quarks.\cite{Laubscher2026PDL_English,Laubscher2026LDP_French}

\subsection{Active surface and coupling to an electron-type \texorpdfstring{\((4,6)\)}{(4,6)}}

The hierarchical closure associated with the proton does not interact with an external \((4,6)\) structure throughout its entire internal volume, but only via a finite ``active surface.'' By \emph{active surface} we designate the minimal envelope of relations that effectively implement the logical coupling with an electron modeled by the \((4,6)\) structure. 

This surface is distinct both from the complete set of internal protonic relations and from the mere aggregate of valence relations. It is defined as the subset of relations that (i) collectively encode the effective charge signature \(+1\), and (ii) remain accessible to the elementary stationary regime associated with the electron.\cite{Laubscher2026LDP_French,Laubscher2026PDL_English}

Within the signed-graph formalism, this coupling is realized through mixed triangles that link vertices associated with the protonic structure to vertices belonging to the \((4,6)\) block.  
Relations constituting the active surface are precisely those that take part in coherent mixed triangles, i.e. triangles which, when combined with the internal sign structure of \((4,6)\), do not induce additional violations of triangular coherence.  
The effective proton charge \(+1\) with respect to the electron is thereby encoded as a global asymmetry that biases the ensemble of accessible mixed triangles toward constructive contributions.\cite{Laubscher2026LDP_French}

For bookkeeping, one may first define a “raw” surface content by weighting the stationary relations of the valence cores with their respective electric charges,
\[
  R_{\text{raw}}
  =
  2\,r_u + r_d
  =
  2\times 276 + 378
  =
  930,
\]
which coincides numerically with the total valence content.
However, not all these relations are equally effective in mediating interactions with the external environment: a subset is primarily dedicated to preserving the internal coherence of the valence cores and does not couple directly to the electron-facing interface.
By incorporating the internal hierarchical structure (locally closed up/down blocks, the relational “sea,” and global closure constraints), one obtains a reduced effective surface comprising on the order of \(R_{\text{surf}} \sim 500\) active relations.\cite{Laubscher2026LDP_French}

A more refined estimate can be formulated in terms of the valence content and an elementary geometrical factor.
Let
\[
  r_{\text{core}} = \frac{r_{\text{val}}}{3}
\]
denote the mean coherence core associated with each valence quark.
By introducing the golden ratio \(\varphi = (1+\sqrt{5})/2\) as a candidate relational optimization parameter, the effective active surface is defined as
\[
  R_{\text{surf}} = \varphi\,r_{\text{core}}
  = \varphi\,\frac{r_{\text{val}}}{3}.
\]
For \(r_{\text{val}}=930\), this yields
\[
  R_{\text{surf}} = \varphi\times 310 \simeq 501.59,
\]
which replaces the crude order-of-magnitude estimate \(R_{\text{surf}}\sim 500\) by an expression formulated exclusively in terms of the discrete valence content and a single geometric factor.\cite{Laubscher2026LDP_French,Laubscher2026PDL_English}

In this framework, the active surface is not to be understood as a geometric boundary in the conventional sense, but rather as a finite relational “skin”: it is defined as the subset of protonic relations that remains available, once internal closure constraints have been satisfied, to participate in the formation of coherent mixed triangles with an external \((4,6)\) structure.
Within this setting, the fine-structure constant is reinterpreted as a measure of the relative efficiency of this coupling, obtained by comparing the proton’s active surface, the minimal electron closure, and the proton–electron mass ratio.
This reinterpretation yields a structural representation of \(\alpha\) as a function of \(\mu\), \(r_{\text{val}}\), and \(\varphi\), which we now derive.

\section{The Fine-Structure Constant Interpreted as a Surface Coherence Fraction}

\subsection{Construction of the active surface}

The proton architecture derived within the PDL framework provides a natural starting point for a structural representation of the fine-structure constant.
The total stationary valence content
\[
  r_{\text{val}} = 930
\]
is distributed over three valence cores, and its mean value per core,
\[
  r_{\text{core}} = \frac{r_{\text{val}}}{3} = 310,
\]
serves as a quantitative measure of a typical local coherence reservoir.
The effective number of surface relations is then obtained by multiplying this mean value by a geometrical optimisation factor, which is identified with the golden ratio \(\varphi = (1+\sqrt{5})/2\),
\[
  R_{\text{surf}} = \varphi\,r_{\text{core}} = \varphi\,\frac{r_{\text{val}}}{3}.
\]
Numerically, this yields
\[
  R_{\text{surf}} \simeq 1.618\times 310 \simeq 501.59,
\]
which refines the earlier approximate estimate \(R_{\text{surf}}\sim 500\) and expresses the active surface solely in terms of the valence content and a single dimensionless factor.\cite{Laubscher2026LDP_French,Laubscher2026PDL_English}

Within this interpretation, the golden ratio is not introduced as an externally imposed “mystical” constant.
Rather, it is regarded as a candidate invariant characterising an optimal redistribution of coherence from a stationary core to its accessible surface, subject to the combined constraints of hierarchical closure and minimal leakage.
The quantity \(R_{\text{surf}}\) therefore enumerates, in purely combinatorial terms, the active channels through which the proton can establish coherent couplings to an external \((4,6)\) structure without compromising its internal stability.

\subsection{PDL expression for the fine-structure constant}

Let \(\mu = m_p/m_e\) denote the proton–electron mass ratio.
Within the PDL framework, the fine-structure constant is reinterpreted as a surface-to-volume coherence fraction, quantifying the ratio between the internal coherence budget of the proton and the residual fraction that remains accessible at its active surface.
A compact, parameter-free expression for the inverse fine-structure constant is then proposed:
\[
  \alpha_{\text{PDL}}^{-1}
  =
  \frac{\mu}{R_{\text{surf}}^2}.
\]
Using the definition of \(R_{\text{surf}}\), this can be recast as
\[
  \alpha_{\text{PDL}}^{-1}
  =
  \frac{\mu}{\varphi^2}\left(\frac{3}{r_{\text{val}}}\right)^2,
\]
which makes explicit that the only continuous input parameter is the mass ratio \(\mu\), whereas \(r_{\text{val}}\) and \(\varphi\) enter, respectively, as a discrete valence parameter and a purely geometrical optimisation factor.\cite{Laubscher2026PDL_English,Laubscher2026LDP_French}

For the numerical values \(r_{\text{val}}=930\), \(\mu\simeq 1836.15\) and \(R_{\text{surf}}\simeq 501.59\), one obtains
\[
  \alpha_{\text{PDL}}^{-1}
  \simeq
  \frac{1836.15}{(501.59)^2}
  \simeq 137.02,
\]
to be compared with the empirical value \(\alpha^{-1}\simeq 137.036\).
The relative deviation is of order a few \(10^{-4}\), and this level of agreement is achieved without introducing any additional continuous free parameter in the construction.\cite{Laubscher2026PDL_English,Laubscher2026LDP_French}
Equivalently, one may write
\[
  \alpha_{\text{PDL}}^{-1}
  =
  \frac{\mu}{\varphi^2}
  \left(\frac{3}{930}\right)^2,
\]
which highlights that the numerical proximity to the observed value results from the combined effect of the discrete valence content, the mass ratio, and the optimisation factor \(\varphi\).

\subsection{Interpretational remarks}

Within this relational framework, the reduced Planck constant \(\hbar\) quantifies the minimal action per coherence cycle of the elementary \((4,6)\) closure associated with the electron, whereas the fine-structure constant \(\alpha\) characterizes the fraction of the proton’s internal coherence budget that is effectively projected onto its active boundary layer.
Accordingly, the product \(\alpha\,\hbar\) can be interpreted as the minimal action available per cycle of surface-mediated coupling between a \((4,6)\) block and the protonic hierarchy, rather than as a bare, unexplained coupling coefficient.\cite{Laubscher2026PDL_English,Laubscher2026LDP_French}

From this standpoint, the dimensionless quantity \(\alpha\) no longer plays the role of an exogenous parameter of the electromagnetic interaction.
Instead, it becomes the quantitative marker of the manner in which a composite stationary regime allocates its finite coherence budget between internal hierarchical closure and external interaction channels.
A substantially larger value of \(\alpha\) would correspond to an overly active surface, thereby undermining the stability of bound states such as the hydrogen atom, whereas a significantly smaller value would render the interface too inert to support stable electronic binding.
The empirically observed magnitude of \(\alpha\) may thus be interpreted as an optimization between preserving the hierarchical closure of the proton and maintaining a minimally sufficient relational “skin” that enables couplings to an electron-type structure.

\section{Effective Coulomb potential and Schr\"odinger description}

\subsection{Standard Coulomb coupling in quantum theory}

In the conventional quantum-mechanical treatment of the hydrogen atom, the proton is idealized as a pointlike static source that generates a Coulomb potential
\[
  V(r)
  =
  -\frac{e^2}{4\pi\varepsilon_0}\,\frac{1}{r}
  =
  -\alpha\,\hbar c\,\frac{1}{r},
\]
where \(\alpha = e^2/(4\pi\varepsilon_0\hbar c)\) denotes the fine-structure constant.
The electron wave function \(\psi(\mathbf{r},t)\) satisfies the nonrelativistic Schr\"odinger equation
\[
  i\hbar \frac{\partial \psi}{\partial t}(\mathbf{r},t)
  =
  \left[
    -\frac{\hbar^2}{2m_e}\Delta
    -\alpha\,\hbar c\,\frac{1}{r}
  \right]\psi(\mathbf{r},t),
\]
and the associated stationary eigenvalue problem reproduces the standard hydrogenic energy spectrum.
For a fixed value of \(\alpha\), the bound-state energies are given by
\[
  E_n
  =
  -\frac{1}{2}m_e c^2 \frac{\alpha^2}{n^2},
  \qquad
  n = 1,2,\dots,
\]
and the corresponding Bohr radius is
\[
  a_0 = \frac{\hbar}{m_e c\,\alpha}.
\]
Within this framework, \(\alpha\) is treated as a fundamental dimensionless parameter that characterizes the strength of the electromagnetic interaction.

\subsection{PDL-based effective Coulomb potential}

Within the PDL framework, the functional form of the Coulomb potential is preserved, while its coupling coefficient is reinterpreted in terms of the proton’s active surface and the proton–electron mass ratio. Substituting \(\alpha\) with \(\alpha_{\text{PDL}}\), as constructed in the preceding section, we introduce an effective potential
\[
  V_{\text{eff}}(r)
  =
  -\alpha_{\text{PDL}}\,\hbar c\,\frac{1}{r},
  \qquad
  \alpha_{\text{PDL}}
  =
  \frac{R_{\text{surf}}^2}{\mu},
\]
where
\[
  R_{\text{surf}} = \varphi\,\frac{r_{\text{val}}}{3},
  \qquad
  \mu = \frac{m_p}{m_e}.
\]
The corresponding time-dependent Schr\"odinger equation takes the form
\[
  i\hbar \frac{\partial \psi}{\partial t}(\mathbf{r},t)
  =
  \left[
    -\frac{\hbar^2}{2m_e}\Delta
    -\alpha_{\text{PDL}}\,\hbar c\,\frac{1}{r}
  \right]\psi(\mathbf{r},t),
\]
which is formally identical to the conventional hydrogenic Schr\"odinger equation, with the crucial distinction that the effective coupling is no longer treated as an independent empirical constant. Instead, it emerges as a derived quantity encoding an internal coherence ratio of the proton–electron system.\cite{Laubscher2026PDL_English,Laubscher2026LDP_French}

\subsection{Hydrogenic spectrum with \texorpdfstring{\(\alpha_{\text{PDL}}\)}{alpha\_PDL}}

Since the radial Schrödinger equation depends solely on the product \(\alpha \hbar c\) entering the Coulomb potential, the standard derivation of the hydrogenic spectrum carries over unchanged upon replacing \(\alpha\) by \(\alpha_{\text{PDL}}\).
Accordingly, the bound-state energy levels are given by
\[
  E_n^{\text{PDL}}
  =
  -\frac{1}{2} m_e c^2 \frac{\alpha_{\text{PDL}}^2}{n^2},
  \qquad
  n = 1,2,\dots,
\]
and the Bohr radius is modified to
\[
  a_0^{\text{PDL}}
  =
  \frac{\hbar}{m_e c\,\alpha_{\text{PDL}}}.
\]
Adopting the numerical estimate \(\alpha_{\text{PDL}}^{-1}\simeq 137.02\), the relative deviation between \(\alpha_{\text{PDL}}\) and the empirical fine-structure constant \(\alpha\) is of order a few \(10^{-4}\). Consequently,
\[
  \frac{E_n^{\text{PDL}} - E_n}{E_n}
  \sim
  2\,\frac{\alpha_{\text{PDL}} - \alpha}{\alpha},
  \qquad
  \frac{a_0^{\text{PDL}} - a_0}{a_0}
  \sim
  -\frac{\alpha_{\text{PDL}} - \alpha}{\alpha},
\]
remain at the same relative magnitude.
Within this level of accuracy, the hydrogen binding energy at the Bohr scale and the corresponding characteristic length scale are correctly reproduced, while the Coulomb coupling attains a structural interpretation in terms of surface coherence and the electron–proton mass ratio.\cite{Laubscher2026PDL_English,Laubscher2026LDP_French}

From this perspective, the PDL framework is not intended as an alternative to the conventional Schrödinger description at atomic length scales.
Instead, it provides a relational interpretation of the coupling parameter in the Coulomb potential by associating \(\alpha\) with the finite number of active surface channels of the proton and with the relative coherence cost of the electron and proton stationary regimes.
The standard hydrogenic equations remain formally unchanged, but the constants that appear in them are no longer treated as purely empirical parameters.

\section{Discussion and outlook}

\subsection{Conceptual status of the construction}

The analysis developed in this work demonstrates a concrete compatibility between the standard Schr\"odinger description of the hydrogen atom and the PDL relational framework. At the phenomenological level, the nonrelativistic Schr\"odinger equation with a Coulomb potential remains intact: the hydrogenic energy spectrum and the Bohr radius preserve their usual expressions once the fine-structure constant is specified. However, the coupling constant entering the Coulomb potential is no longer treated as a primitive parameter; instead, it is reinterpreted as an emergent ratio characterizing internal coherence, formulated in terms of the proton’s active surface, the proton–electron mass ratio, and the discrete valence content.\cite{Laubscher2026PDL_English,Laubscher2026LDP_French}

From a conceptual perspective, this construction indicates that several quantities traditionally regarded as fundamental constants may admit a structural origin within a pre-metric relational framework. The minimal \((4,6)\) closure associated with the electron determines the granularity of action through \(\hbar\), while the proton’s internal organization gives rise to a finite number of effective surface channels and to a natural scale for a surface-to-volume coherence fraction. In this setting, the fine-structure constant emerges as the quantitative trace of how a composite stationary regime allocates its finite coherence budget between internal closure and external coupling, rather than as a purely exogenous input parameter.\cite{Laubscher2026IntroPDL,Laubscher2026PDL_English}

\subsection{Limitations and perspectives}

Several limitations of the present approach must be highlighted.

First, the numerical agreement \(\alpha_{\text{PDL}}^{-1}\simeq 137.02\) versus \(\alpha^{-1}\simeq 137.036\) is obtained within a simplified combinatorial representation of the proton and depends critically on the specific identification of \(r_{\text{val}}=930\) and \(R_{\text{sea}}=10\,087\) as optimal solutions of the PDL optimisation functional. A more exhaustive combinatorial exploration is required to establish the uniqueness and stability of this particular architecture, and to determine whether alternative integer assignments could yield comparable or possibly superior agreement.\cite{Laubscher2026PDL_English,Laubscher2026LDP_French}

Second, the Schr\"odinger equation has been adopted as a fundamental postulate. No attempt has been made to derive nonrelativistic quantum dynamics from PDL axioms, nor to reconstruct the full formalism of quantum field theory. Consequently, the scope of the present work is deliberately restricted: it offers a structural reinterpretation of a single coupling constant within an otherwise standard quantum-mechanical framework, rather than a comprehensive relational reformulation of quantum theory. Extending the PDL framework to incorporate dynamical laws, multi-particle sectors, and quantum statistics remains an open research programme.\cite{Laubscher2026IntroPDL,Laubscher2026PDL_English}

Third, the same PDL architecture has previously been employed to propose candidate expressions for Boltzmann’s constant and an effective gravitational coupling, which reproduce the correct orders of magnitude but still exhibit discrepancies at the level of a few percent or more. These constructions invoke additional elements, such as a small coherence-leakage parameter and hierarchical amplification factors, whose detailed combinatorial origin is not yet fully clarified. A systematic refinement of these ingredients, together with a more stringent comparison with general relativity and equilibrium as well as nonequilibrium thermodynamics, is necessary to assess the broader physical viability of the framework.\cite{Laubscher2026PDL_English,Laubscher2026LDP_French}

Notwithstanding these limitations, the present compatibility result suggests that a nontrivial subset of atomic physics can be reformulated within a purely relational, graph-theoretic description of microscopic structure. Future work will aim to sharpen the combinatorial analysis of nucleon architectures, investigate potential corrections to the effective Coulomb potential at short distances, generalise the framework to other bound systems such as helium-like ions, and examine whether characteristic PDL-induced deviations from standard quantum predictions could be probed by high-precision spectroscopic measurements or scattering experiments.\cite{Laubscher2026PDL_English,Laubscher2026LDP_French}

\section{Conclusion}

In this work, we have analysed the consistency between the standard nonrelativistic Schr\"odinger description of the hydrogen atom and the Projective Dynamic Logo (PDL) framework.
Beginning from a purely relational formulation on signed graphs, in which the minimal \((4,6)\) structure is construed as a logical prototype of an electron at rest and the proton is encoded as a two-level hierarchical architecture with three valence cores and an associated relational sea, we have derived a structural expression for the fine-structure constant.

The central ingredients of the construction are the discrete valence content \(r_{\text{val}}=930\), the proton–electron mass ratio \(\mu=m_p/m_e\), and an effective active surface
\[
  R_{\text{surf}} = \varphi\,\frac{r_{\text{val}}}{3},
\]
where \(\varphi\) is interpreted as a relational optimisation factor.
These quantities combine into a parameter-free expression
\[
  \alpha_{\text{PDL}}^{-1}
  =
  \frac{\mu}{R_{\text{surf}}^2}
  =
  \frac{\mu}{\varphi^2}\left(\frac{3}{r_{\text{val}}}\right)^2,
\]
which yields \(\alpha_{\text{PDL}}^{-1}\simeq 137.02\), in close agreement with the empirical value \(\alpha^{-1}\simeq 137.036\), with a relative deviation of order a few \(10^{-4}\).\cite{Laubscher2026PDL_English,Laubscher2026LDP_French}

Upon inserting \(\alpha_{\text{PDL}}\) into the conventional Coulomb potential,
\[
  V_{\text{eff}}(r) = -\alpha_{\text{PDL}}\hbar c\,\frac{1}{r},
\]
we recover the standard hydrogenic spectrum and the Bohr radius, with corrections controlled by the small discrepancy between \(\alpha_{\text{PDL}}\) and \(\alpha\).
Consequently, the usual Schr\"odinger phenomenology of hydrogen is preserved, while the Coulomb coupling constant attains a structural interpretation in terms of surface coherence and the mass ratio within the PDL architecture.\cite{Laubscher2026PDL_English,Laubscher2026LDP_French}

The present analysis does not yet provide a complete derivation of quantum dynamics from the PDL axioms, nor does it encompass higher-order corrections or multi-particle correlations.
It nonetheless demonstrates that a strictly relational framework, formulated with minimal assumptions, can reproduce in a quantitatively nontrivial manner the numerical value of a central coupling constant of atomic physics.
Further research will be necessary to extend this relational reconstruction to additional constants and interactions, and to identify potential experimental signatures capable of discriminating PDL-based structures from more conventional microscopic models.\cite{Laubscher2026IntroPDL,Laubscher2026PDL_English}

\bibliographystyle{unsrt} % ou autre
\bibliography{references}

\end{document}
